% ─────────────────────────────────────────────────────────────────────────────
%  Sección 4.X  Arquitectura del Módulo de Gestión de Ejercicios
%  Insertar en el Capítulo 4: Diseño, después del módulo de Autenticación
%  y antes del módulo OCR/NLP.
% ─────────────────────────────────────────────────────────────────────────────

\subsection{Arquitectura del Módulo de Gestión de Ejercicios}

El módulo de gestión de ejercicios administra el ciclo de vida completo
de las actividades de escritura dentro de la plataforma. Su diseño contempla
dos actores con responsabilidades distintas: el tutor, que selecciona
ejercicios del catálogo y los asigna a sus alumnos con una vigencia
determinada; y el alumno, que consulta los ejercicios disponibles
clasificados por su estado. El módulo también incorpora un proceso
automatizado que cierra los ejercicios cuya fecha de vencimiento ha
sido alcanzada, garantizando la coherencia del estado del sistema sin
intervención manual.

La arquitectura del módulo se divide en tres componentes principales:

\begin{enumerate}
    \item \textbf{Capa de Presentación (Front-End):} Compuesta por los
    componentes \texttt{TabEjercicios.jsx} (vista del tutor para gestionar
    asignaciones), \texttt{TabProximos.jsx}, \texttt{TabCompletados.jsx}
    y \texttt{TabVencidos.jsx} (vistas del alumno que clasifican los
    ejercicios según su estado). Cada componente consulta su endpoint
    correspondiente y renderiza la información sin lógica de negocio
    propia.

    \item \textbf{Capa de Lógica de Negocio (Back-End):} Implementada
    en el controlador \texttt{ejercicios.py} mediante FastAPI. Gestiona
    la consulta del catálogo, la activación y desactivación de ejercicios
    por parte del tutor, y la consulta de ejercicios en sus tres estados
    posibles desde la perspectiva del alumno. Adicionalmente, el proceso
    \texttt{cerrar\_ejercicio.py} implementa un scheduler que verifica
    y cierra automáticamente los ejercicios vencidos cada hora.

    \item \textbf{Capa de Persistencia (Base de Datos):} El módulo opera
    sobre dos tablas principales: \texttt{Ejercicios}, que actúa como
    catálogo inmutable de actividades disponibles en la plataforma, y
    \texttt{Ejercicios\_Tutor}, que representa la asignación de un
    ejercicio del catálogo a un tutor específico, con su propio ciclo de
    vida controlado por el campo \texttt{id\_estatus} y la fecha
    \texttt{fecha\_desactivacion}.
\end{enumerate}

% ─────────────────────────────────────────────────────────
\subsubsection{Diseño del Ciclo de Vida de un Ejercicio}

Un ejercicio atraviesa los siguientes estados a lo largo de su ciclo de
vida dentro del sistema:

\begin{enumerate}
    \item \textbf{Disponible en catálogo:} el ejercicio existe en la tabla
    \texttt{Ejercicios} con \texttt{id\_estatus = 1} y puede ser
    seleccionado por cualquier tutor para su asignación.

    \item \textbf{Activo:} el tutor ha creado un registro en
    \texttt{Ejercicios\_Tutor} con \texttt{id\_estatus = 1}. A partir de
    este momento, el ejercicio es visible para los alumnos del tutor en
    su panel de ejercicios próximos.

    \item \textbf{Vencido:} el ejercicio transiciona a
    \texttt{id\_estatus = 2} por uno de dos mecanismos: desactivación
    manual por parte del tutor mediante el endpoint de eliminación, o
    cierre automático por el scheduler cuando la fecha
    \texttt{fecha\_desactivacion} es anterior a la fecha y hora actuales
    del servidor. Un ejercicio vencido deja de aparecer en la lista de
    próximos del alumno y pasa a la lista de vencidos si no fue entregado,
    o permanece en completados si el alumno ya realizó un intento.
\end{enumerate}

\noindent \textit{Nota de diseño}: el sistema no elimina físicamente los
registros de \texttt{Ejercicios\_Tutor}. La transición entre estados se
gestiona exclusivamente mediante el campo \texttt{id\_estatus}, lo que
preserva la trazabilidad histórica de las asignaciones y permite al tutor
consultar ejercicios previamente activos sin pérdida de información.

% ─────────────────────────────────────────────────────────
\subsubsection{Diseño del Flujo de Interacción}

El módulo contempla tres flujos principales de interacción.

\textbf{Flujo de Asignación de Ejercicio (Tutor):}
\begin{enumerate}
    \item El tutor accede a su panel y solicita el catálogo de ejercicios
    disponibles mediante una petición \texttt{GET} al endpoint público
    \texttt{/api/ejercicios/}.
    \item El tutor selecciona un ejercicio y establece opcionalmente una
    fecha de vencimiento. El Front-End envía una petición \texttt{POST}
    al endpoint \texttt{/api/ejercicios/tutor} con el identificador del
    ejercicio y la fecha de cierre.
    \item El Back-End verifica mediante JWT que el \texttt{id\_usuario}
    del token coincide con el \texttt{id\_usuario} del payload, evitando
    que un tutor active ejercicios en nombre de otro.
    \item Se verifica que el tutor no tenga ya activo el mismo ejercicio.
    En caso de duplicado, se retorna \texttt{400 Bad Request}.
    \item Se inserta un nuevo registro en \texttt{Ejercicios\_Tutor} con
    \texttt{id\_estatus = 1}. El ejercicio queda inmediatamente disponible
    para los alumnos del tutor.
\end{enumerate}

\textbf{Flujo de Consulta de Ejercicios (Alumno):}

Desde la perspectiva del alumno, los ejercicios se presentan clasificados
en tres pestañas. La lógica de clasificación reside íntegramente en el
Back-End mediante tres consultas SQL independientes:

\begin{itemize}
    \item \textbf{Próximos:} ejercicios activos (\texttt{id\_estatus = 1})
    asignados por el tutor del alumno que aún no tienen un intento
    registrado por ese alumno. Se excluyen mediante una subconsulta
    \texttt{NOT IN} sobre la tabla \texttt{Intentos}.

    \item \textbf{Completados:} ejercicios que tienen al menos un intento
    registrado por el alumno. Se retorna únicamente el intento más reciente
    por ejercicio, determinado mediante la función de ventana
    \texttt{ROW\_NUMBER() OVER (PARTITION BY id\_ejercicio\_tutor ORDER BY
    fecha\_envio DESC)}, incluyendo la calificación y retroalimentación
    del tutor si ya fueron registradas.

    \item \textbf{Vencidos:} ejercicios con \texttt{id\_estatus = 2} que
    no tienen ningún intento registrado por el alumno, es decir, los que
    el alumno no alcanzó a entregar antes del cierre.
\end{itemize}

\paragraph{Flujo de Cierre Automático (Scheduler):}
\begin{enumerate}
    \item Al iniciar el servidor, el evento \texttt{startup} de FastAPI
    invoca \texttt{iniciar\_scheduler()}, que registra una tarea periódica
    con una frecuencia de ejecución de una hora mediante la biblioteca
    APScheduler.
    \item Cada hora, la tarea \texttt{cerrar\_ejercicios\_vencidos()}
    ejecuta un \texttt{UPDATE} sobre la tabla \texttt{Ejercicios\_Tutor},
    transitando a \texttt{id\_estatus = 2} todos los registros cuya
    \texttt{fecha\_desactivacion} sea distinta de nulo y anterior a la
    fecha y hora actuales del servidor (\texttt{GETDATE()}).
    \item La operación es idempotente: ejecutarla múltiples veces sobre
    los mismos datos no produce efectos secundarios, ya que la condición
    \texttt{id\_estatus = 1} del filtro impide modificar registros ya
    cerrados.
\end{enumerate}

% ─────────────────────────────────────────────────────────
\subsubsection{Diseño de APIs}

Para soportar los flujos descritos, el diseño de la API contempla los
siguientes endpoints en el controlador \texttt{ejercicios.py}:

\textbf{Endpoint 1:} Consulta del Catálogo

El endpoint \texttt{/api/ejercicios/} opera mediante el método \texttt{GET}
y no requiere autenticación. Retorna todos los ejercicios con
\texttt{id\_estatus = 1} disponibles en la tabla \texttt{Ejercicios}.

\textbf{Endpoint 2:} Consulta de Ejercicios Activos del Tutor

El endpoint \texttt{/api/ejercicios/tutor/\{id\_usuario\}} opera mediante
el método \texttt{GET} y requiere autenticación JWT con rol de tutor.
Retorna los ejercicios que el tutor ha activado y que se encuentran
actualmente con \texttt{id\_estatus = 1}.

\textbf{Endpoint 3:} Activación de Ejercicio

El endpoint \texttt{/api/ejercicios/tutor} opera mediante el método
\texttt{POST} y requiere autenticación JWT con rol de tutor. Recibe el
identificador del ejercicio del catálogo y una fecha de vencimiento
opcional.

\begin{lstlisting}[language=json, caption={Payload de solicitud del endpoint de activación de ejercicio}, label={lst:ejercicio-create}]
/* Solicitud */
{
  "id_ejercicio": <integer>,
  "id_usuario":   <integer>,
  "fecha_fin":    "datetime | null"
}

/* Respuesta 201 Created */
{
  "message": "Ejercicio activado correctamente"
}
\end{lstlisting}

\textbf{Endpoint 4:} Desactivación Manual de Ejercicio

El endpoint \texttt{/api/ejercicios/tutor/\{id\_usuario\}/\{id\_ejercicio\}}
opera mediante el método \texttt{DELETE} y requiere autenticación JWT con
rol de tutor. Actualiza el campo \texttt{id\_estatus} a 2 en el registro
correspondiente de \texttt{Ejercicios\_Tutor}, sin eliminar el registro
físicamente.

\textbf{Endpoints 5, 6 y 7:} Consulta de Ejercicios del Alumno

Los endpoints \texttt{/api/ejercicios/alumno/\{id\_alumno\}/proximos},
\texttt{/api/ejercicios/alumno/\{id\_alumno\}/completados} y
\texttt{/api/ejercicios/alumno/\{id\_alumno\}/vencidos} operan mediante
el método \texttt{GET} y requieren autenticación JWT con rol de alumno.
En los tres casos, el Back-End verifica que el \texttt{id\_alumno} de la
ruta coincida con el \texttt{id\_usuario} del token, retornando
\texttt{403 Forbidden} en caso de discrepancia.

La Tabla~\ref{tab:endpoints-ejercicios} resume los endpoints del módulo
con su método y nivel de acceso correspondiente.

\begin{table}[H]
\centering
\caption{Endpoints del módulo de gestión de ejercicios}
\label{tab:endpoints-ejercicios}
\begin{tabular}{|l|l|l|}
\hline
\textbf{Endpoint} & \textbf{Método} & \textbf{Acceso requerido} \\
\hline
\texttt{/api/ejercicios/}
    & GET    & Público \\
\texttt{/api/ejercicios/tutor/\{id\_usuario\}}
    & GET    & Tutor (JWT) \\
\texttt{/api/ejercicios/tutor}
    & POST   & Tutor (JWT) \\
\texttt{/api/ejercicios/tutor/\{id\_usuario\}/\{id\_ejercicio\}}
    & DELETE & Tutor (JWT) \\
\texttt{/api/ejercicios/alumno/\{id\_alumno\}/proximos}
    & GET    & Alumno (JWT) \\
\texttt{/api/ejercicios/alumno/\{id\_alumno\}/completados}
    & GET    & Alumno (JWT) \\
\texttt{/api/ejercicios/alumno/\{id\_alumno\}/vencidos}
    & GET    & Alumno (JWT) \\
\hline
\end{tabular}
\end{table}